\Askhsh{14190}{4}
\begin{ylh}{2}
    \item 2.1. Οι Πράξεις και οι Ιδιότητές τους
    \item 4.2. Ανισώσεις 2ου Βαθμού
    \item 6.1. Η Έννοια της Συνάρτησης
    \item 6.2. Γραφική Παράσταση Συνάρτησης
\end{ylh}
Δίνεται η συνάρτηση $f(x) = x^{2} + x + 1\ ,\ x \in R$.

\begin{alist}
\item Να αποδείξετε ότι η γραφική παράσταση της συνάρτησης $f\ $βρίσκεται πάνω από τον άξονα $x΄x$.
\monades{8}
\item Να αποδείξετε ότι για οποιονδήποτε πραγματικό αριθμό $a \neq - \dfrac{1}{2}$, τα σημεία της γραφικής παράστασης της $f$ με τετμημένες $a$ και $–\ a\ –\ 1$ έχουν την ίδια τεταγμένη.
\monades{8}
\item Θεωρούμε μεταβλητό σημείο Μ της γραφικής παράστασης της $f$ με τετμημένη $\beta > 0$. Από το Μ φέρνουμε παράλληλες ευθείες προς τους άξονες $x΄x$ και $y΄y$ και έστω Α και Δ τα σημεία τομής αυτών των ευθειών με τους άξονες, όπου το Α ανήκει στον $x΄x$ και το Δ στον $y΄y$. Αποδείξτε ότι η περίμετρος του ορθογωνίου ΟΑΜΔ είναι $\left\lbrack \sqrt{2}(\beta + 1) \right\rbrack^{2}$.

\monades{9}
\end{alist}

